\documentclass[journal,twoside,web]{ieeecolor}
\usepackage{tmi}
\usepackage{cite}
\usepackage{amsmath,amssymb,amsfonts}
\usepackage{algorithmic}
\usepackage{graphicx}
\usepackage{textcomp}
\usepackage{booktabs}
\usepackage[normalem]{ulem}
\usepackage{colortbl}
\usepackage{hyperref}
\hypersetup{
    colorlinks,
    citecolor=blue,
    filecolor=black,
    linkcolor=black,
    urlcolor=blue
}
\usepackage{subcaption}
\captionsetup[subfigure]{singlelinecheck=false, skip=0pt}

% Remove later
\usepackage{color,soul}
\newcommand{\TG}[1]{\color{blue}\textsc{From Tristan: }#1\color{black}}
\newcommand{\MD}[1]{\color{magenta}\textsc{From Mathieu: }#1\color{black}}
\newcommand{\HL}[1]{\hl{#1}}
% End of remove

\def\BibTeX{{\rm B\kern-.05em{\sc i\kern-.025em b}\kern-.08em
    T\kern-.1667em\lower.7ex\hbox{E}\kern-.125emX}}
\markboth{\journalname, VOL. XX, NO. XX, XXXX 2020}
{Author \MakeLowercase{\textit{et al.}}: Preparation of Papers for IEEE TRANSACTIONS ON MEDICAL IMAGING}
\begin{document}
\title{An Evaluation of Mixed Precision Arithmetic in ANTs Registration}
\author{Mathieu Dugr\'e , Yohan Chatelain, and Tristan Glatard
\thanks{This work was supported in part by a Natural Sciences and Engineering Research Council of Canada (NSERC) Postgraduate Scholarship-Doctoral (PGS D) scholarship awarded to M. Dugr\'e. This work was also supported by Compute Canada and the Digital Research Alliance of Canada who provided the compute resources for our experiments.}
\thanks{M. Dugr\'e and T. Glatard are with the Department of Computer Science and Software Engineering, Concordia University, Montreal, Quebec, Canada (e-mail: research@mathdugre.me).}
\thanks{Y. Chatelain and T. Glatard are with the Centre for Addiction and Mental Health (CAMH), Toronto, Ontario, Canada (e-mail: {yohan.chatelain@gmail.com, tristan.glatard@camh.ca})}}

\maketitle

\begin{abstract}
	Image registration is key to MRI analysis, allowing comparison between two
	or more images. Its high compute cost limits analysis of larger cohorts and
	time sensitive applications. Mixed precision is a technique leveraging high-
	and low-precision arithmetic to accelerate computations while maintaining
	accuracy. Its success in accelerating deep learning models led to a broad
	adoption, but its adoption remains limited in other scientific domains due
	to the limited understanding of its effect in complex application. In this
	work, we propose to study the effects of mixed precision in ANTs
	registration---a common MRI registration pipeline. We use VPREC to simulate
	arbitrary precision arithmetic, and develop a model to predict the minimal
	precision required for convergence per iteration of the optimization
	process. Our results suggest that using low precision (FP32 \& FP24) is
	sufficient in ANTs registration. Our model tends to underestimate the number
	of precision bits required by 11.6\%, on average \TG{transition with
	previous sentence is missing} (threshold: $10^{-3}$ \TG{remove mention of
	threshold in abstract, it's almost impossible to understand what it means}).
	Our work advances the understanding of mixed precision in ANTs registration,
	by characterizing when the method \TG{``the method'' is ambiguous} is
	beneficial. Therefore, laying the ground for practical acceleration of the
	pipeline \TG{This sentence isn't grammatically correct}. Furthermore, our
	theoretical model offers a simple approach to optimizing the precision in
	ANTs registration, offering a building block for more accurate predictions
	\TG{why to more accurate predictions? This is quite ambiguous}.
\end{abstract}

\begin{IEEEkeywords}
	mixed precision, image registration, performance, mri, neuroimaging, ANTs, ITK
\end{IEEEkeywords}

\section{Introduction}
\label{sec:introduction}
Image registration transforms a pair or collection of images to a common space,
enabling comparison between those images. It is a critical step in magnetic
resonance imaging (MRI) processing pipelines as it enable the comparison of
images acquired at different times, from different subjects, or from different
modalities. Therefore, allowing for longitudinal studies, group analysis, or
multi-modal imaging \TG{this sentence isn't grammatically correct}. By
consequence \TG{Frenchicism?}, registration is an important step in toolkits such as the Advanced
Normalization Tools (ANTs) registration framework~\cite{Avants2020-xx}---a
widely adopted tool in the neuroimaging community.

MRI registration incurs a significant computational cost, often accounting for a majority of the runtime for MRI processing pipelines. As a result, optimizing registration performance is crucial for improving the overall efficiency of MRI processing---enabling studies with larger cohorts and clinical applications requiring processing in timely manner. Our previous study~\cite{Dugre2025-fs} identified two primary performance bottlenecks in MRI processing pipelines: (1) interpolation computation in the registration process and (2) memory bandwidth limitations.

Mixed precision arithmetic leverages a mix of high- and low-precision data
formats to accelerate applications while maintaining acceptable accuracy. Data
formats with smaller bit width have faster arithmetic operations \TG{Data
formats don't "have" operations} and reduce the memory footprint. Therefore, it
\TG{ambiguous pronoun reference: it refers to mixed precision while the writing
flow suggests it refers to data formats} has the potential to alleviate both
bottlenecks. However, while mixed precision has the potential for significant
performance gains, it is unclear when and where it can be applied in complex
scientific pipelines like medical image registration without compromising result
integrity \TG{integrity refers to something else, use accuracy instead}.

While mixed precision has the potential for significant performance gains, determining the appropriate precision either spatially, temporally, or both, is challenging in complex or large-scale applications. We conduct a detailed analysis of mixed precision formats in the context of ANTs registration pipeline~\cite{Avants2020-xx}. This work contributes to the understanding and application of mixed precision formats in MRI registration by doing the following:
\begin{itemize}
	\item \textbf{Empirical Evaluation:} Demonstrate that low-precision formats (i.e., FP32 and FP24) satisfy the numerical requirements of ANTs registration, safely relaxing the default reliance on resource-intensive FP64 arithmetic.
	\item \textbf{Theoretical Bounds:} Introduce a predictive model establishing the minimum computational precision required at each optimization iteration to guarantee convergence.
	\item \textbf{Heuristic Validation:} Validate these mathematical bounds against simulation benchmarks, providing a reliable heuristic to guide dynamic precision tuning for future performance optimization.
\end{itemize}

We organize the remainder of this paper as follows. Section~\ref{sec:background}
introduces the concept used to build our methods and discusses related work on
mixed precision computation and numerical sensitivity of MRI registration.
Section~\ref{sec:theory} details the theoretical analysis of minimum precision
requirements, adapting the theorem by Hayford et al.~\cite{Hayford2024-kb}.
Section~\ref{sec:experiments} describes our experimental setup, including the
ANTs registration pipeline, our use of VPREC to simulate arbitrary data formats,
and the dataset used. Finally, we report and discuss our results in
Section~\ref{sec:results} and Section~\ref{sec:discussion}, respectively.
Section~\ref{sec:conclusion} concludes the paper.

\section{Background}
\label{sec:background}

\subsection{Reduced and Mixed Precision Arithmetic}
The computational complexity of arithmetic operations is approximately quadratic relative to the bit-width of the data format utilized~\cite{Wang2019-qe}. Consequently, enabling applications to utilize reduced bit-width formats offers significant performance benefits. Traditionally, pipeline developers have relied on \textit{binary64} (FP64) and \textit{binary32} (FP32) to balance resource requirements with numerical accuracy. However, many applications may not necessitate the high dynamic range or precision of these formats and could be significantly accelerated through mixed precision strategies.

Mixed precision optimizes performance by leveraging a combination of high- and low-bit-width arithmetic. These variations in precision can be applied either temporally (across iterations) or spatially (across specific code sections). While mixed precision promises higher performance gains than static mixed precision methods, its implementation in complex pipelines is often technically challenging, typically requiring extensive and intrusive code modifications.

\subsection{Prior Work on Mixed Precision}
Mixed precision has been widely adopted in deep learning (DL), where the escalating parameters and computational demands of modern models have driven the adoption of lower precision formats, such as half-precision (FP16), bfloat16 (BF16), and sub-8-bit representations. This significantly reduces memory bandwidth usage while accelerating matrix multiplications. The foundation of modern mixed precision training was established by Micikevicius et al.~\cite{Micikevicius2017-sf}, who introduced critical techniques to prevent accuracy degradation when training with FP16, including maintaining an FP32 master copy of weights, loss scaling, and high-precision accumulation. PyTorch~\cite{Paszke2019-vk} and TensorFlow~\cite{Abadi2016-vm}, popular DL frameworks, provide native support for Automatic Mixed Precision (AMP) utilizing static heuristics mapping operations to predefined precisions.

To further mitigate the scaling complexities inherent in standard IEEE formats, specialized hardware representations have emerged, such as Brain Floating Point (BF16)~\cite{Kalamkar2019-br} and Nvidia's TensorFloat-32 (TF32)~\cite{Kharya2020-th}, as well as 8-bit and 4-bit formats representing the frontier for large language model applications~\cite{Sevegnani2025-zd}.

Despite these advancements in DL, the adoption of mixed precision in classical scientific pipelines, such as neuroimaging, remains limited. These scientific pipelines generally lack ``out-of-the-box'' support for mixed precision and often default to FP64 to ensure stability. Standard implementations of AMP fail to encapsulate diverse data distributions and optimization topologies. Alternative solutions encompass static analysis---like DAISY~\cite{Darulova2018-wv} and Herbie~\cite{Panchekha2015-au}---which guarantee error bounds offline but output overly conservative configurations and are challeging to apply to complex application. Dynamic analysis directly tracks real-time program execution to allocate precision, seen in approaches like SERP~\cite{Rios2021-it}, CoMP~\cite{Dai2025-ek}, and DQMQ~\cite{Wang2023-iu}, but these techniques lack a theoretical foundation and often incurs large computational overhead. 

To alleviate the heavy runtime exploration costs of dynamic tuning, theoretical bounds characterizing computation precision requirements present a promising avenue~\cite{Sakr2017-ec, Tu2023-os, Wang2024-re}. Providing a rigorously generalizable mathematical anchor, Hayford et al.~\cite{Hayford2024-kb} linked gradient-based optimization convergence directly to computational precision thresholds based on parameter norm, the Lipschitz constant, and loss gradient norm variables. This theoretical framework offers the foundation for tuning precision at runtime without the need for costly dynamic exploration, but has yet to be expanded for precision tuning or applied to complex scientific pipelines such as medical image registration.

\subsection{Simulation of Arbitrary Precision}
Commodity hardware supports only a narrow subset of data formats---typically among FP64, FP32, FP16, BF16, \textit{Float24} (FP24), and TF32. Developing hardware standards for new formats is resource-intensive and acquisition of new hardware is costly. Therefore, it is essential to determine which formats best balance performance and accuracy for a specific pipeline.

While theoretical analysis can determine the precision requirements for a sequence of arithmetic operations, the cumulative complexity of such analysis is often infeasible for complex or large-scale pipelines.

Alternatively, dynamic analysis evaluates the precision-accuracy trade-off by varying precision across multiple executions. Tools like the VPREC backend of Verificarlo~\cite{Denis2015-zf,Chatelain2019-vw} enable the runtime simulation of arbitrary precision, enabling dynamic analysis. VPREC is non-intrusive, requiring only the recompilation of source code to instrument a pipeline. However, VPREC simulation is computationally expensive, incurring overheads ranging from $2.6\times$ to $16.8\times$~\cite{Chatelain2019-vw}. Thus, while VPREC is a powerful exploratory tool for defining numerical requirements, it is not suitable for production-level optimization.

\subsection{Numerical Sensitivity in MRI Registration}
MRI registration is commonly implemented as an iterative optimization process
that aligns a moving image to a fixed template. Within the ANTs framework, this
is performed via a multi-stage pipeline each operating on a multi-resolution
pyramid. Numerical errors introduced during the initial coarse alignment can
propagate and amplify during subsequent high-resolution stages.

A study by Mirhakimi et al.~\cite{Mirhakimi2025-qb} demonstrated that numerical
uncertainty in ANTs linear registration can result in mean displacement
artifacts of up to 0.2mm---a magnitude comparable to typical motion observed
during MRI acquisitions. Furthermore, the choice of similarity metric influences
stability; Mutual Information (MI) typically exhibits higher sensitivity to
numerical noise than Cross-Correlation (CC). These findings underscore that
numerical stability in registration is not merely a matter of metric convergence
but also of anatomical integrity \TG{not sure what anatomical integrity means. I
would rather underscore that numerical stability is data dependent, anatomy
being only one aspect of the data}. Misalignments or topological distortions
introduced by numerical inaccuracies can lead to erroneous clinical
interpretations. Therefore, balancing accuracy and performance in MRI
registration is critical to ensure reliable outcomes.

Accuracy thresholds in this context are highly application-specific and remain
poorly defined. Determining these thresholds would require extensive empirical
validation across diverse clinical applications. While establishing these limits
is critical to fully realizing the potential of mixed precision, it remains a
significant challenge in the field and is outside the primary scope of this
work. \TG{I find this paragraph misplaced and ambiguous. I would rather just
omit it, or move it to results.}

\section{Theoretical Analysis}
\label{sec:theory}
To analyze the impact of rounding errors on convergence, we adapted the theoretical framework proposed by Hayford et al.~\cite{Hayford2024-kb}. The model characterizes the optimization convergence under mixed precision arithmetic, specifically focusing on the point where rounding noise dominates the gradient signal. While originally derived to bound rounding errors in DL optimization using BF16 arithmetic, the underlying theorem is applicable to any gradient-based optimization, including the iterative registration stages in ANTs.

Let $\theta$ represent the optimization parameters, $\mathcal{L}(\theta)$
the loss function, and $\eta$ the learning rate. In a mixed precision setting
with $p$ bits of mantissa precision, the gradient descent update is perturbed by
an error of $\delta\theta$:
\begin{equation}
	\theta \leftarrow \theta - \eta \nabla\mathcal{L}(\theta + \delta\theta)
\end{equation}
where the rounding error satisfies the relative bound $\|\delta\theta\| \leq 2^{-p} \|\theta\|$.

Under the assumptions that $\mathcal{L}$ is convex and $\nabla \mathcal{L}$ is Lipschitz continuous with constant $K$ such that $K < 1/\eta$:
\begin{equation}
	\|\nabla \mathcal{L}(x) - \nabla \mathcal{L}(y)\| \leq K \|x - y\| \quad \forall x, y,
\end{equation}
the optimization is guaranteed to progress toward the minimum only as \TG{``only
as'' is ambiguous, do you mean "if and only if"?} long as the gradient signal
$\|\nabla \mathcal{L}(\theta)\|$ exceeds the rounding noise. The theorem shows
that if the model \TG{which model? Do you mean ``the optimization''} does not
converge to the minimum, then it will reach a critical zone \TG{``the model''
doesn't reach a zone. Do you mean that the optimal value theta star will?}
bounding the loss gradient norm as follows:
\begin{equation}
	\|\nabla \mathcal{L}(\theta)\| \leq \frac{4+3\sqrt{2}}{2^{p}} K \|\theta\|.
\end{equation}
The constant $4+3\sqrt{2}$ is a symbolic bound derived from the quadratic
inequality required to guarantee $\mathcal{L}(\theta_{t+1}) <
\mathcal{L}(\theta_t)$. By taking the contrapositive, we establish a sufficient
condition for convergence. To ensure the gradient signal is not masked by
precision-induced noise, the precision $p$ must satisfy \TG{I think the logic is
wrong here, it should be ``if the bound is verified then the optimization will
converge''. But if the bound is not verified then the optimization may also
converge, which you should note.}:
\begin{equation}
	p > p_{min} = \log_2 \left( \frac{(4+3\sqrt{2}) K \|\theta\|}{\|\nabla\mathcal{L}(\theta)\|} \right)
	\label{eq:pmin}
\end{equation}
where $p_{min}$ is the minimum number of mantissa bits required at a given
iteration. We modified ITK, the optimization library used by ANTs, to probe $K$,
$\|\nabla \mathcal{L}\|$, and $\|\theta\|$ at each iteration. As in the original
study \TG{not sure what the original study is, be more explicit}, $K$ was estimated using a first-order finite difference:
\begin{equation}
	K \approx \frac{\|\nabla \mathcal{L}(\theta_{t+1}) - \nabla \mathcal{L}(\theta_t)\|}{\|\theta_{t+1} - \theta_t\|}
\end{equation}
By computing Equation~\ref{eq:pmin} dynamically, we determine the minimum
precision per iterations necessary to maintain the accuracy of the ANTs
registration optimization.

The model could be used to dynamically adjust precision at runtime in ANTs, or other gradient-based optimization (e.g., deep learning training), by instrumenting the code to compute $p_{min}$ at each iteration and adjusting the precision accordingly. However, this dynamic adjustment would require extensive code modification and is outside the scope of this work. We focus on evaluating the accuracy of the model in predicting precision requirements.

\section{Experiments}
\label{sec:experiments}
\subsection{ANTs Registration SyN}
The ANTs registration framework~\cite{Avants2020-xx} is a multi-stage, multi-resolution iterative optimization pipeline that aligns a moving image to a fixed template. The default configuration consists of three sequential stages---rigid, affine, and symmetric normalization (SyN)~\cite{Avants2008-ea}---each operating across four resolution levels. These levels use a shrink factors of $8\times4\times2\times1$ with corresponding Gaussian smoothing kernels of $3\times2\times1\times0$, respectively. Prior to optimization, the images are aligned via an initial intensity-based center-of-mass transformation.

By default, the rigid and affine stages utilize mutual information (MI) as the objective
function, while the SyN stage uses cross-correlation (CC). Images are smoothed
prior to down-sampling at each level. The process terminates when the
optimization reaches the maximum number of iterations or when the objective
function converges within a threshold of $10^{-6}$ over the last 10 iterations.
ANTs leverages the registration framework from the Insight Toolkit
(ITK)~\cite{Yoo2002-ve, McCormick2014-ok} for computation.

\subsection{Mixed Precision using VPREC}
Although ANTs provides a flag to enable single precision (FP32) computation, our previous study~\cite{Dugre2025-fs} identified a bug in ITK\footnote{https://github.com/InsightSoftwareConsortium/ITK/issues/4593} where interpolation operations execute in double precision (FP64) regardless of compilation or runtime settings. This results in casting overhead that degrades performance compared to the default FP64 pipeline execution.

To explore the impact of mixed precision, we used Verificarlo with the VPREC backend to simulate specific bit-widths via compile-time instrumentation. We evaluated combinations of exponent ranges (6--8 bits) and mantissa precision (6--23 bits) against the default FP64 baseline.

We hypothesized that precision requirements vary by stage: failures in early alignment may propagate, while high-resolution stages likely necessitate finer precision to ensure convergence. Consequently, we tested precision reduction across various stage combinations as detailed in Table~\ref{tab:vprec_combinations}. We recompiled ANTs and ITK with Verificarlo for runtime simulation of precision reduction. Instrumentation was disabled for metric computation modules, as low-precision noise in these components caused the pipeline to crash due to values deviating from the expected convergence range.

\begin{table}[t]
	\centering
	\begin{tabular}{@{}ccccc@{}}
		\toprule
		\multicolumn{1}{c}{Experiment ID} & \begin{tabular}[c]{@{}c@{}}Initialization\\ (Moving)\end{tabular} & Rigid & Affine & SyN \\ \midrule
		1111                              & 1                                                                 & 1     & 1      & 1   \\
		0000                              & 0                                                                 & 0     & 0      & 0   \\
		0001                              & 0                                                                 & 0     & 0      & 1   \\
		0011                              & 0                                                                 & 0     & 1      & 1   \\
		0111                              & 0                                                                 & 1     & 1      & 1   \\
		0110                              & 0                                                                 & 1     & 1      & 0   \\
		0100                              & 0                                                                 & 1     & 0      & 0   \\
		\bottomrule
	\end{tabular}
	\caption{Configuration of Stage-Specific Precision Reduction. Experiment IDs are encoded as a four-digit binary sequence where 0 represents default FP64 precision and 1 represents VPREC-instrumented mixed precision. The digits correspond to the following pipeline stages: (1) initialization, (2) rigid, (3) affine, and (4) SyN.}
	\label{tab:vprec_combinations}
\end{table}

\subsection{Empirical and Estimated Computation}
The empirical $p_{min}$ was estimated by evaluating VPREC results against the FP64 baseline across relative difference thresholds ($10^{-2}$ to $10^{-5}$). For each threshold, we averaged results across all subjects and identified the minimum bit-width $p$ that maintained the final iteration error within the specified tolerance.

\subsection{Dataset and Infrastructure}
We used the \textit{BMB\_1} site from the \textit{CoRR}
dataset~\cite{Zuo2014-ub} ($N=50$; mean age: $30.8$ [19.9--59.7]; $52\%$
female). Anatomical $T1$-weighted volumes were skull-stripped via the
antsBrainExtraction.sh script~\cite{Avants2020-xx,ANTs2023-du} and
quality-controlled against the \textit{tpl-MNI152NLin2009cAsym}
template~\cite{Ciric2022-mb}. Experiments were executed on the \textit{Rorqual}
cluster of the Digital Research Alliance of Canada using containerized Apptainer
images with architecture-specific instructions disabled to ensure
reproducibility across CPU microarchitectures. \TG{you could link to your github repo here}

\section{Results}
\label{sec:results}

\subsection{FP32 Provides Sufficient Accuracy for ANTs Registration}
We evaluated the impact of mixed precision by measuring the relative difference in similarity metrics compared to the FP64 baseline. Fig.~\ref{fig:ants-space-search} illustrates these differences across combinations of exponent range and mantissa precision.

To improve readability, we omit redundant configurations for the affine and SyN stages from the visualization. We observed that the numerical behavior within a specific stage remained consistent across different configurations, provided the stage itself was instrumented with the same precision parameters.
\begin{figure*}[htb!]
	\centering
	\begin{subfigure}[t]{\textwidth}
		\caption{[01xx] Rigid (MI)}
		\label{subfig:ants-search-01xx}
		\includegraphics[width=\textwidth]{figures/0111-rigid-metric_rel_diff-mean.pdf}
	\end{subfigure}
	\begin{subfigure}[t]{\textwidth}
		\caption{[001x] Affine (MI)}
		\label{subfig:ants-search-001x}
		\includegraphics[width=\textwidth]{figures/0011-affine-metric_rel_diff-mean.pdf}
	\end{subfigure}
	\begin{subfigure}[t]{\textwidth}
		\caption{[0001] SyN (CC)}
		\label{subfig:ants-search-0001}
		\includegraphics[width=\textwidth]{figures/0001-syn-metric_rel_diff-mean.pdf}
	\end{subfigure}
	\caption{ Mean relative difference in similarity metrics compared to the
		FP64 baseline across varied exponent and mantissa configurations
		($N=50$). Subfigures (a)--(c) represent the rigid, affine, and SyN
		stages, respectively, without instrumentation in preceding stages.
		Heatmap colors indicate performance relative to FP64: red denotes
		degradation, green denotes improvement, and white indicates numerical
		equivalence, with the color intensity representing the magnitude of the
		relative difference. Formats with hardware support are framed:
		Float32 (r8-p23), Float24 (r7-p16), TensorFloat32 (r8-p10), and BFloat16
		(r8-p7). The 4-bit binary ID encodes instrumentation (1), default
		precision (0), or not executed (x) for: (1) initialization, (2) rigid,
		(3) affine, and (4) SyN stages.}
	\label{fig:ants-space-search}
\end{figure*}

Our results indicate that low-precision formats, such as FP32 (r8-p23) and FP24
(r7-p16), maintain high accuracy, with relative differences remaining
below $10^{-5}$ compared to FP64. To achieve a relative difference threshold of
$10^{-4}$, the rigid and affine stages require a minimum of 11 bits of
precision, while the SyN stage necessitates at least 13 bits. These requirements
are significantly lower than the 53 bits provided by FP64.

\subsection{Current Stage Precision Primarily Dictates Final Registration Accuracy}
Table~\ref{tab:stage-accuracy} illustrates how precision at different stages of
the pipeline affects the final similarity metric. We specifically examined the
BF16 format to observe how errors propagate from the initial rigid alignment
through to the final SyN stage.

When the rigid, affine, and SyN stages were all instrumented in BF16 (Exp. ID
\uline{011}\textbf{1}), the accuracy degraded progressively, resulting in a
final mean relative difference of $8.68\%$. However, when the SyN stage was
reverted to FP64 while maintaining BF16 for all prior stages (Exp. ID
\uline{011}\textbf{0}), the accuracy recovered significantly. In this
configuration, the final similarity metric was $0.54\%$ better than the FP64
baseline, despite the low-precision initialization.

This ``accuracy recovery'' suggests that the numerical errors introduced in
early stages are not additive; instead, the final registration accuracy is
governed by the precision of the SyN stage rather than the cumulative error
propagated from preceding steps.

\TG{conclude on the possibility to use BF16 alltogether, if only the output of
the SyN stage is to be used}

% Exception for the initialization step
\definecolor{RedHL}{RGB}{253, 104, 100}
\definecolor{GreenHL}{RGB}{103, 253, 154}
\begin{table}[htb!]
	\centering
	\begin{tabular}{@{}llll@{}}
		\toprule
		\multicolumn{1}{c}{\textbf{Exp. ID}} & \multicolumn{1}{c}{\textbf{Rigid (MI)}} & \multicolumn{1}{c}{\textbf{Affine (MI)}} & \multicolumn{1}{c}{\textbf{SyN (CC)}} \\ \midrule
		{\uline{011}\textbf{1}}              & \cellcolor{RedHL}{2.34e-02}             & \cellcolor{RedHL}{1.86e-01}              & \cellcolor{RedHL}{8.68e-02}           \\
		{\uline{011}\textbf{0}}              & \cellcolor{RedHL}{2.34e-02}             & \cellcolor{RedHL}{1.86e-01}              & \cellcolor{GreenHL}{-5.38e-03}        \\ \bottomrule
	\end{tabular}
	\caption{Impact of Stage-Specific Instrumentation on Final Registration Accuracy (BF16 vs. FP64). The experiment IDs indicate whether the current and prior stages are instrumented (1) or maintained at default precision (0)---only the SyN stage is varied in these experiments. Red and green highlights indicate worse and better similarity relative to FP64, respectively.}
	\label{tab:stage-accuracy}
\end{table}

\subsection{Predicted $p_{min}$ Yields Conservative Estimates for Convergence}
We compared the predicted $p_{min}$ via Equation~\ref{eq:pmin} with empirical
VPREC estimations. The empirical $p_{min}$ represents the minimum precision
required to maintain the mean relative metric difference below specific
thresholds ($10^{-2}$ to $10^{-5}$).

As shown in Fig.~\ref{fig:ants-pmin-overview}, the predicted $p_{min}$ (dashed
line) generally acts as a conservative upper bound for the empirical thresholds,
indicating that the predicted number of bits is sufficient for convergence. In
the rigid and affine stages, $p_{min}$ typically remains below the 23 bits
provided by FP32, while the requirements for the SyN stage are lower---typically
below 15 bits (satisfied by FP24). However, we observe that the model
underestimates the precision required for SyN levels 1 and 2.
\begin{figure*}[htb!]
	\centering
	\includegraphics[width=\textwidth]{figures/pmin-threshold-last-iteration/pmin-overview.pdf}
	\caption{
		Comparison of theoretically predicted $p_{min}$ (dashed line) and empirical VPREC estimations (solid lines). Values represent the mean precision bits required across subjects ($N=50$) at each iteration. VPREC curves indicate the precision necessary to maintain relative difference thresholds ($T$) from $10^{-2}$ to $10^{-4}$. Metrics $U$ (Under) and $O$ (Over) denote the area ratio percentages where the theoretical model underestimates or overestimates the empirical requirement, respectively. Horizontal light gray lines indicate precision bits for standard floating-point hardware formats.}
	\label{fig:ants-pmin-overview}
\end{figure*}

Table~\ref{tab:area-ratio-summary} summarizes the percentage of the area where our prediction is above (overestimation) or below (underestimation) the VPREC thresholds across all stages. For a relative difference of $10^{-3}$ ($0.1\%$), our model underestimates the required precision by $11.57\%$ on average. Conversely, it overestimates by $34.57\%$, providing a safety margin that ensures the convergence guarantee in the majority of stage-level combinations. We observed a large standard deviation across these values, indicating that our model accuracy varies significantly across specific stage-level combinations.
\begin{table}[htb!]
	\centering
	\begin{tabular}{@{}lll@{}}
		\toprule
		\multicolumn{1}{c}{\textbf{Threshold}} & \multicolumn{1}{c}{\textbf{Under}} & \multicolumn{1}{c}{\textbf{Over}} \\ \midrule
		$10^{-2}$                              & $6.15\pm6.90$                      & $50.59\pm43.58$                   \\
		$5\times10^{-3}$                       & $7.41\pm7.03$                      & $43.32\pm35.24$                   \\
		$10^{-3}$                              & $11.57\pm9.75$                     & $34.57\pm30.70$                   \\
		$10^{-4}$                              & $16.79\pm11.40$                    & $26.03\pm22.61$                   \\ \bottomrule
	\end{tabular}
	\caption{Summary of Prediction Accuracy via Area Ratio Analysis. Values represent the mean and standard deviation $(\mu \pm \sigma)$ of the area where the theoretical $p_{min}$ under or over estimates the empirical VPREC thresholds across all stage-level combinations. A smaller "Over" ratio signifies a lower overhead, but a tighter margin for convergence.}
	\label{tab:area-ratio-summary}
\end{table}

\section{Discussion}
\label{sec:discussion}
Our evaluation of mixed precision within the ANTs framework reveals that standard low-precision formats (FP32 and FP24) are generally sufficient for robust image registration. Furthermore, our theoretical model offers an effective heuristic for bounding dynamic precision requirements. The following sections discuss the implications, limitations, and practical integration of these findings.

\subsection{Numerical Stability and Convergence Floors}
Our results establish the first ``numerical floor'' for the ANTs and ITK frameworks. While low-precision formats such as Float24 (r7-p16) maintain accuracy, reducing the representation to 6 bits or fewer in either the exponent or mantissa consistently led to execution instabilities, including segmentation faults or indefinite hangs. Given the complex, multi-branched architecture of the ITK-based optimization pipeline, we attribute these failures to numerical noise triggering divergent execution paths or the failure to satisfy convergence thresholds, leading to unexpected execution states.

Furthermore, the ``catastrophic failures'' observed during low-precision initialization (Exp. 1111) underscore the sensitivity of the initial global alignment. Because initialization defines the search space for all subsequent stages, rounding errors at this stage are particularly detrimental; at very low mantissa precision, these errors resulted in severe topological discontinuities and loss of anatomical structures in the transformed volumes. This suggests that while later stages exhibit a degree of resilience to numerical noise---even recovering accuracy lost in previous iterations---the center-of-mass alignment must be preserved at higher precision to ensure a stable registration.

\subsection{Theoretical Assumptions and Model Predictions}
The theoretical model adapted from Hayford et al.~\cite{Hayford2024-kb} assumes a convex loss function. However, in medical image registration with high-resolution 3D MRI volumes and complex similarity metrics like cross-correlation, the cost function is likely non-convex. Despite this theoretical limitation, our empirical results demonstrate that our $p_{min}$ model (Eq.~\ref{eq:pmin}) remains an effective heuristic.

Prior work in deep learning suggest that non-convexity may not invalidate convergence. Gur-Ari et al.~\cite{Gur-Ari2018-pm} demonstrated that stochastic gradient descent confines optimization to a low-dimensional subspace spanned by the top Hessian eigenvectors, which can exhibit locally convex properties, despite a globally non-convex loss landscape. Additionally, Liu et al.~\cite{Liu2020-pe} showed that over-parametrized neural networks often satisfy the Polyak-Lojasiewicz~\cite{Lojasiewicz1963-xd,Polyak1963-nd} condition, enabling convergence to a global minimum despite non-convex landscapes despite non-convex loss landscape.

The ANTs registration problem differs from deep learning optimization in a critical way: transformation parameters directly affect anatomical integrity \TG{not sure what this means, rephrase} by modifying spatial voxel mappings, making over-parametrization arguments inapplicable. However, several factors may create a locally convex optimization subspace: the multi-resolution pyramid progressively refines the search space, Gaussian smoothing regularizes the loss landscape, and explicit regularization constraints further stabilize optimization. Future work should characterize the ANTs loss landscape to validate these theoretical assumptions.

Our model attempts to capture complex optimization dynamics via heuristic. Its predictions are conservative, underestimating the bit requirements by only 6--17\% across difference thresholds and tends to over-predict by 26--50\%. Over-prediction preserves the convergence guarantee but limits the potential speedup from lower-bit formats, while under-prediction risks convergence issues. Balancing these factors is essential for maximizing performance without compromising numerical stability.

\subsection{Implementation of Mixed Precision}
Our current experiments incur a large computational cost by simulating each range-precision combination with VPREC. Utilizing our theoretical model to predict $p_{min}$---and subsequently modifying the precision during runtime---would allow for dynamic precision tuning without repeated executions. This capability shifts precision optimization from a costly offline simulation to a dynamic, runtime process, enabling temporal (per iteration) or spatial (per code section) mixed precision in research environments.

However, fine-grained mixed precision presents two primary technical challenges. First, frequent precision casting can cause overhead that outscales the arithmetic benefits. Second, implementing such tuning is difficult in large, legacy codebases like ITK. In our methods the similarity metric calculation is computed in FP64 to preserve stability. In practice, this would create a large casting overhead at each iteration and be challenging to implement spatial mixed precision in ITK. Future research must align the similarity metric's precision with the optimization pipeline to simplify configuration and eliminate casting bottlenecks.

\subsection{Limitations and Clinical Relevance}
\TG{I would avoid using the word ``clinical'' here, clinical validation would involve a whole other range of data and metrics}

While our model effectively predicts the minimum precision requirement, several components limit its generalizability and require further research. First, our current model makes precision predictions based on the FP64 loss landscape. Dynamically adjusting the computational precision based on these predictions would modify the loss landscape, potentially affecting the accuracy of subsequent predictions. Second, we utilized the $L_2$ norm consistently with Hayford et al.~\cite{Hayford2024-kb}, which does not account for the varying units and scales of affine transformation parameters. Future studies should investigate scale-invariant norms, such as the Mahalanobis norm. Third, since previous work demonstrated that similarity metrics significantly influence numerical stability~\cite{Mirhakimi2025-qb}, the interaction between other metric types and mixed precision requires deeper investigation.

Finally, the accuracy thresholds in this study were strictly defined via relative numerical differences from the FP64 baseline. These pure mathematical thresholds lack direct clinical interpretation. To confidently transition mixed precision registration into clinical deployment, future work must evaluate the impact of minor registration differences on concrete downstream tasks. This involves ensuring reliability in clinical applications such as volume estimation in voxel-based morphometry, longitudinal tracking of lesions, or high-stakes surgical planning, where the margin for topological distortion is critical.

\section{Conclusion}
\label{sec:conclusion}
In this work, we evaluated the feasibility and impact of using mixed precision arithmetic in the ANTs registration pipeline. Our empirical analysis demonstrates that standard single-precision (FP32) and even lower-precision formats like FP24 are sufficient to maintain registration accuracy comparable to the default double-precision (FP64) baseline. We observed that while the initialization stage is highly sensitive to numerical precision, subsequent optimization stages exhibit a degree of resilience, allowing for ``accuracy recovery'' where high precision in the final stage compensates for lower precision in earlier steps.

Furthermore, we introduced a theoretical model to predict the minimum precision ($p_{min}$) required for convergence at each iteration. While our model tends to be conservative---often overestimating the required bit-width---it provides a safety margin that ensures convergence. This theoretical framework offers a promising path toward dynamic mixed precision registration, where precision can be tuned on-the-fly to improve performance without compromising results accuracy.

Overall, our findings suggest that the stringent requirement for FP64 in neuroimaging pipelines may be relaxed. By carefully selecting precision formats for specific stages or iterations, we can unlock significant performance gains on modern hardware. Future work should address technical implementation challenges to translate these findings into practical high-performance tools for large-scale neuroimaging.

\bibliography{paper.bib}
\bibliographystyle{IEEEtran}
\end{document}
